% !TEX output_directory = ../publica
\documentclass[11pt, a4paper]{article}

\usepackage[utf8]{inputenc}
\usepackage[T1]{fontenc}
\usepackage[portuguese]{babel}
\usepackage[sfdefault]{roboto}
% Geometria exata baseada nas marcações milimétricas (Área de texto começando em 1.5cm)
\usepackage[top=1.5cm, bottom=1.5cm, left=1cm, right=1cm, headheight=12pt, headsep=0.4cm]{geometry}
\usepackage{graphicx}

% --- MINI-CABEÇALHO E RODAPÉ EM TODAS AS PÁGINAS ---
\usepackage{fancyhdr}
\pagestyle{fancy}
\fancyhf{} 
\renewcommand{\headrulewidth}{0pt} 
% Mini-cabeçalho mais claro (gray!50), alinhado ao topo, com BIMESTRE no final
\fancyhead[R]{\scriptsize\bfseries\textcolor{gray!50}{\MakeUppercase{MATEMÁTICA} $\vert$ Bloco SEMANA 10 $\vert$ 2º BIMESTRE}}
\fancyfoot[C]{\footnotesize Página \thepage} 

\usepackage{setspace}
\setstretch{1.15}
\usepackage{parskip}
\setlength{\parskip}{2pt}
\setlength{\parindent}{0pt}

\usepackage{amsmath, amsfonts, amssymb}
\usepackage{nccmath}
\usepackage{mathastext}
\usepackage{xcolor}
\definecolor{CorTitUm}{HTML}{FF6F61}
\definecolor{CorTitDois}{HTML}{FF6F61}
\definecolor{CorTitTres}{HTML}{658894}
\definecolor{CorLinha}{HTML}{B0B0B0}   

% --- BOX RESPONDA (ESTILO BLOCO DE CITAÇÃO C/ LINHA LATERAL) ---
\usepackage[most]{tcolorbox}
\newtcolorbox{boxresponda}{
    colback=white,         
    colframe=gray!60,      
    boxrule=0pt,           
    leftrule=3pt,          
    sharp corners,         
    left=8pt, right=0pt,   
    top=2pt, bottom=2pt,
    fontupper=\small      
}

\usepackage{titlesec}
\titleformat{\section}{\color{CorTitUm}\fontsize{17}{20}\bfseries}{\thesection}{0em}{}
\titlespacing*{\section}{0pt}{5pt}{6pt} 
\titleformat{\subsection}{\color{black}\fontsize{13}{15}\bfseries\raggedright}{\thesubsection}{0em}{}
\titlespacing*{\subsection}{0pt}{1pt}{5pt} 

\usepackage{multicol}
\setlength{\columnsep}{1cm}
\setlength{\columnseprule}{0.5pt}
\def\columnseprulecolor{\color{CorLinha}}

\usepackage{adjustbox}
\usepackage{tikz}
\usepackage{pgfplots}
\usepackage{circuitikz}
\usetikzlibrary{shadows, shapes.misc, positioning, calc}
\pgfplotsset{compat=1.18}

\begin{document}

% --- CABEÇALHO PRINCIPAL ESTILO CADERNO (Apenas 1ª página) ---
\noindent
\begin{minipage}[t]{0.12\linewidth}
    \vspace{0pt}
    % Substitua 'logo.png' pelo brasão da sua escola se desejar compilar com logo.
    % \includegraphics[width=\linewidth, height=1.65cm, keepaspectratio]{logo.png} 
\end{minipage}
\hfill
\begin{minipage}[t]{0.85\linewidth}
    \vspace{0pt}
    % BLOCO 1 (Canto inferior esquerdo reto - southwest)
    \begin{tcolorbox}[colback=white, colframe=gray!50, arc=4pt, sharp corners=southwest, boxrule=0.8pt, left=10pt, right=10pt, top=6pt, bottom=6pt]
        \begin{minipage}[c]{0.65\linewidth}
            \Large\bfseries CADERNO DE ATIVIDADES
        \end{minipage}
        \hfill
        \begin{minipage}[c]{0.33\linewidth}
            \raggedleft\small\bfseries\textcolor{gray!70!black}{\MakeUppercase{MATEMÁTICA}}
        \end{minipage}
    \end{tcolorbox}
    \vspace{0.15cm} % Distanciamento entre os blocos
    % BLOCO 2 (Canto superior esquerdo reto - northwest)
    \begin{tcolorbox}[colback=white, colframe=gray!50, arc=4pt, sharp corners=northwest, boxrule=0.8pt, left=10pt, right=10pt, top=2pt, bottom=2pt]
        \small\bfseries\textcolor{gray!85}{Ano/Série:} \textcolor{black}{3º ANO / PRÉ-VESTIBULAR} \hspace{2.5cm} \textcolor{gray!85}{Bloco:} \textcolor{black}{SEMANA 10}
    \end{tcolorbox}
\end{minipage}

\vspace{0.7cm} % Empurra a linha grossa exatamente para a marca de 3.9cm
\noindent\rule{\linewidth}{1.2pt} 
\vspace{0.4cm}
% --- FIM DO CABEÇALHO ---

\raggedcolumns
\begin{multicols*}{2}

\vspace{0.5cm}
\subsection*{Semana 10: O Círculo e a Circunferência}
\vspace{-0.2cm}\noindent\rule{\linewidth}{1.2pt} 
\vspace{0.2cm}

\textbf{1.} (ENEM 2012) Uma pizzaria oferece duas opções de tamanho: a pizza média, com diâmetro de \textbf{30 cm}, e a pizza grande, com diâmetro de \textbf{40 cm}. Um cliente deseja saber a diferença real de tamanho (área) entre as duas opções antes de decidir o pedido. 

Considerando o valor de $\pi = 3$, a diferença de área entre a pizza grande e a pizza média, em centímetros quadrados, é exatamente de:

a) 75\\[0.3cm]
b) 100\\[0.3cm]
c) 300\\[0.3cm]
d) 525\\[0.3cm]
e) 2.100\\[0.3cm]

\vspace{0.4cm}
\noindent\textcolor{CorLinha}{\rule{\linewidth}{0.5pt}} 
\vspace{0.2cm}

\textbf{2.} (ENEM - Adaptada) O proprietário de um sítio deseja cercar um lago de formato perfeitamente circular utilizando uma tela especial de arame. Ele verificou que a distância do centro do lago até a sua margem é de \textbf{15 m}. Sabe-se que a tela é vendida exclusivamente em rolos fechados de \textbf{50 m}. Considere $\pi = 3,14$.

Com base nestes dados, o proprietário concluiu que precisaria comprar \textbf{3 rolos} de tela para dar uma volta completa no lago e ainda guardar uma sobra de \textbf{55,8 m}. 

Refute a conclusão do proprietário, provando matematicamente por que a quantidade de rolos e a sobra calculadas por ele estão incorretas. Determine a quantidade real de rolos necessários.

\vspace{2.5cm}

\vspace{0.4cm}
\noindent\textcolor{CorLinha}{\rule{\linewidth}{0.5pt}} 
\vspace{0.2cm}

\textbf{3.} (FUVEST) Em um parque, há uma grande praça circular com raio de \textbf{20 m}. Bem no centro dessa praça, foi construído um chafariz, também de formato circular, com raio de \textbf{5 m}. A prefeitura pretende pavimentar com blocos de concreto toda a área da praça, deixando intacta apenas a área já ocupada pelo chafariz. 

Determine, em metros quadrados e em função de $\pi$, a área exata que receberá a pavimentação de concreto.

\vspace{2.5cm}

\vspace{0.4cm}
\noindent\textcolor{CorLinha}{\rule{\linewidth}{0.5pt}} 
\vspace{0.2cm}

\textbf{4.} (ENEM 2014) Uma bicicleta de marcha única possui o pneu dianteiro com \textbf{60 cm} de diâmetro. Um ciclista amador vai realizar um percurso em linha reta medindo exatos \textbf{1.884 m}. 

Considerando a aproximação $\pi = 3,14$, determine o número total de voltas completas que o pneu dianteiro dessa bicicleta dará até o final do trajeto percorrido.

a) 100\\[0.3cm]
b) 500\\[0.3cm]
c) 1.000\\[0.3cm]
d) 10.000\\[0.3cm]
e) 100.000\\[0.3cm]

\vspace{0.4cm}
\noindent\textcolor{CorLinha}{\rule{\linewidth}{0.5pt}} 
\vspace{0.2cm}

\textbf{5.} (Estilo ENEM) Um radar de trânsito fixo consegue monitorar uma área da rodovia que possui formato geométrico de um setor circular. O manual do equipamento informa que ele possui um raio de alcance de \textbf{12 km} e sua lente monitora um ângulo de varredura de exatos \textbf{60 graus}.

Para calcular a área monitorada, um estudante de engenharia apresentou a seguinte resolução em seu relatório:

\textit{"Se a área do círculo total é calculada por $\pi \cdot R^2$, temos $\pi \cdot 12^2 = 144\pi$. Como o ângulo é de 60 graus, e o círculo tem 360 graus, eu devo dividir a área total por 3, já que $3 \times 60 = 180$, que é a metade da circunferência. Portanto, a área de varredura é de $48\pi$ km²."}

\begin{boxresponda}
\textbf{Responda:} Qual foi o erro conceitual primário na linha de raciocínio do estudante ao estabelecer a proporção do setor circular? Corrija o cálculo estruturando a proporção correta e apresente a área real coberta pelo radar.
\end{boxresponda}

\vspace{0.4cm}
\noindent\textcolor{CorLinha}{\rule{\linewidth}{0.5pt}} 
\vspace{0.2cm}

\textbf{6.} (ENEM 2018) Um projeto paisagístico baseia-se em um canteiro central de formato quadrado, com lado medindo \textbf{10 m}. O arquiteto responsável inscreveu dentro deste quadrado o maior espelho d'água circular possível, de modo que a borda do círculo tangencie os quatro lados do canteiro. O espaço do quadrado que não foi ocupado pela água será preenchido com grama natural.

\begin{center}
\begin{adjustbox}{max width=\linewidth}
\begin{tikzpicture}
\definecolor{asfalto}{HTML}{2F3640}
\definecolor{grama}{HTML}{27AE60}
\definecolor{agua}{HTML}{3498DB}

% --- APLICANDO DROP SHADOW E ROUNDED CORNERS ---
% LAYER 1: Canteiro quadrado com grama e sombra
\fill[drop shadow={opacity=0.2, shadow xshift=2pt, shadow yshift=-2pt}, grama, rounded corners=2pt] (-2.5,-2.5) rectangle (2.5,2.5);

% LAYER 2: Espelho d'água (Círculo)
\fill[agua, opacity=0.85] (0,0) circle (2.5);
\draw[white, thick] (0,0) circle (2.5);

% LAYER 3: Cotas com máscara de legibilidade
\draw[|<->|, thick, asfalto] (-2.5,-2.9) -- (2.5,-2.9) node[midway, fill=white, inner sep=2pt] {\textbf{10 m}};
\draw[dashed, thick, white] (0,0) -- (2.5,0);
\fill[asfalto] (0,0) circle (2pt);
\node[fill=agua, inner sep=1pt, text=white] at (1.25, 0.25) {\textbf{R}};

\end{tikzpicture}
\end{adjustbox}
\end{center}

Considerando $\pi = 3,1$, a área aproximada que receberá a aplicação da grama natural, em metros quadrados, é igual a:

a) 22,5\\[0.3cm]
b) 40,0\\[0.3cm]
c) 69,0\\[0.3cm]
d) 78,5\\[0.3cm]
e) 100,0\\[0.3cm]

\vspace{0.4cm}
\noindent\textcolor{CorLinha}{\rule{\linewidth}{0.5pt}} 
\vspace{0.2cm}

\textbf{7.} (FUVEST) Analise detalhadamente a afirmação a seguir sobre as propriedades fundamentais da geometria plana:

\textit{"Ao se duplicar o comprimento do raio de um círculo qualquer, o valor do seu perímetro dobra, porém o valor da sua área quadruplica, uma vez que a dimensão radial se encontra elevada ao quadrado na determinação da superfície."}

Julgue a afirmativa como VERDADEIRA ou FALSA. Justifique a sua resposta apresentando o equacionamento algébrico que comprova o comportamento da área e do perímetro após a alteração.

\vspace{2.5cm}

\vspace{0.4cm}
\noindent\textcolor{CorLinha}{\rule{\linewidth}{0.5pt}} 
\vspace{0.2cm}

\textbf{8.} (ENEM 2013) Uma pista de atletismo foi projetada na forma de um grande retângulo ladeado por dois semicírculos idênticos em suas extremidades. O trecho retangular central tem \textbf{100 m} de comprimento. A largura total da pista, que corresponde exatamente ao diâmetro de cada semicírculo, é de \textbf{60 m}.

\begin{center}
\begin{adjustbox}{max width=\linewidth}
\begin{tikzpicture}
\definecolor{pista}{HTML}{E67E22}
\definecolor{linha}{HTML}{FFFFFF}
\definecolor{grama}{HTML}{2ECC71}

% --- APLICANDO DROP SHADOW E LAYERS ---
% LAYER 1: Sombra e Fundo Verde (Grama do Estádio)
\fill[drop shadow={opacity=0.2, shadow xshift=3pt, shadow yshift=-3pt}, grama, rounded corners=8pt] (-4,-2.5) rectangle (4,2.5);

% LAYER 2: Pista de Atletismo (Retângulo + 2 Semicírculos)
\fill[pista] (-2,-1.5) rectangle (2,1.5);
\fill[pista] (-2,0) circle (1.5);
\fill[pista] (2,0) circle (1.5);

% LAYER 3: Área central em grama
\fill[grama] (-2,-1) rectangle (2,1);
\fill[grama] (-2,0) circle (1);
\fill[grama] (2,0) circle (1);

% LAYER 4: Linhas de marcação
\draw[linha, thick, dashed] (-2,-1.25) rectangle (2,1.25);
\draw[linha, thick, dashed] (-2,0) circle (1.25);
\draw[linha, thick, dashed] (2,0) circle (1.25);
\draw[linha, thick] (-2,-1) -- (-2,1);
\draw[linha, thick] (2,-1) -- (2,1);

% LAYER 5: Cotas com fundo máscara (inner sep)
\draw[|<->|, thick] (-2, 1.8) -- (2, 1.8) node[midway, fill=grama, inner sep=2pt, text=white] {\textbf{100 m}};
\draw[|<->|, thick] (3.8, -1.5) -- (3.8, 1.5) node[midway, fill=grama, inner sep=2pt, text=white] {\textbf{60 m}};

\end{tikzpicture}
\end{adjustbox}
\end{center}

Um atleta de corrida treina exatamente sobre a marcação da linha mais interna desta pista. Considerando o valor de $\pi = 3$, qual é a distância exata percorrida pelo atleta ao completar uma volta inteira na pista?

a) 260 m\\[0.3cm]
b) 320 m\\[0.3cm]
c) 380 m\\[0.3cm]
d) 400 m\\[0.3cm]
e) 480 m\\[0.3cm]

\vspace{0.4cm}
\noindent\textcolor{CorLinha}{\rule{\linewidth}{0.5pt}} 
\vspace{0.2cm}

\textbf{9.} (ENEM 2020) No mecanismo interno de uma máquina industrial, engrenagens menores transmitem rotação para engrenagens maiores de forma acoplada. O sistema possui uma engrenagem motora de raio \textbf{10 cm} firmemente acoplada a uma engrenagem movida de raio \textbf{25 cm}. 

Considerando o princípio das grandezas proporcionais em roldanas e polias, determine a quantidade de voltas completas que a engrenagem maior dará quando a engrenagem menor completar exatamente \textbf{10 voltas}. Apresente a memória de cálculo.

\vspace{2.5cm}

\vspace{0.4cm}
\noindent\textcolor{CorLinha}{\rule{\linewidth}{0.5pt}} 
\vspace{0.2cm}

\textbf{10.} (FUVEST - Adaptada) O comprimento de um arco é determinado pelo ângulo central que o "enxerga". Considere uma circunferência maciça cujo raio mede exatamente \textbf{18 cm}. Um recorte específico desta circunferência forma um arco cujo ângulo central mede \textbf{40 graus}. 

O comprimento exato da linha correspondente a esse arco, expresso em função de $\pi$, é de:

a) $2\pi$ cm\\[0.3cm]
b) $4\pi$ cm\\[0.3cm]
c) $6\pi$ cm\\[0.3cm]
d) $8\pi$ cm\\[0.3cm]
e) $12\pi$ cm\\[0.3cm]

\end{multicols*}
\end{document}